\section{Controle de Versão e Integração Contínua}
\label{sec:controle-versao-integracao-continua}

A gestão de configuração de software compreende a identificação, o controle e o registro das
alterações realizadas nos artefatos de um projeto ao longo do tempo, de modo que seja possível
reconstituir versões anteriores e identificar a origem de cada mudança \cite{sommerville2011}. O
controle de versão é o mecanismo que concretiza esse registro sobre o código-fonte: cada
alteração é capturada em um commit, com autor, data e descrição, e o histórico resultante
permite recuperar estados anteriores do sistema.

O módulo utilizou Git \cite{git2026} para esse fim, com ramificações (branches) por
funcionalidade, o que isolou o desenvolvimento do módulo das demais partes do Banco de
Especialistas até que as alterações fossem verificadas. Os repositórios foram hospedados no
GitLab \cite{gitlab2026}, adotado como padrão pelo Laboratório Assert para os projetos do Polo
de Inovação. A integração contínua automatiza a execução de verificações, como a suíte de
testes Pytest, a cada commit enviado ao repositório; neste projeto, essa execução é realizada
por uma pipeline, impedindo que alterações que quebrem as regras de negócio sejam incorporadas à
base principal.

% TO-DO: Imagem de uma pipeline de integração contínua no GitLab.

Cabe observar que o controle de versão atua sobre o código-fonte, e não sobre os dados da
aplicação. A rastreabilidade das alterações cadastrais dos especialistas, tratada na seção \ref{sec:auditoria-rastreabilidade-versionamento},
é um problema distinto, resolvido no domínio da aplicação por meio de solicitações,
snapshots e registro de eventos.
